%--------------------------------------------------------------------------------
%
%
%--------------------------------------------------------------------------------
\pagebreak
\section*{SUB — Integer Subtraction}
\addcontentsline{toc}{section}{SUB - Integer Subtraction}

\instructionentry{Syntax}{
  \texttt{SUB RegR, RegB, RegA} \\
  \texttt{SUB RegR, RegB, SignedImmed15}
  
  \texttt{SUB[.D/W/H/B] RegR, SignedImmed13( RegB )} \\
  \texttt{SUB[.D/W/H/B] RegR, RegX ( RegB )}
}

\instructionentry{Format}{

	The SUB instruction uses the register, the immediate, the indexed and the register indexed instruction formats. \\
	
	\begin{flushleft}	

	\drawInstrFormat
  	{0,1,3,5,6.5,8.5,9.5,11.5,16}
  	{\OpCodeGrpALU,\OpCodeSUB,RegR, 0, RegB, 0, RegA, 0}
	
	\drawInstrFormat
  	{0,1,3,5,6.5,8.5,16}
  	{\OpCodeGrpALU,\OpCodeSUB, RegR, 1, RegB, val}

	\drawInstrFormat
  	{0,1,3,5,6.5,8.5,9.5,16}
  	{\OpCodeGrpMEM,\OpCodeSUB, RegR, 0, RegB, dw, ofs}
 
	\drawInstrFormat
  	{0,1,3,5,6.5,8.5,9.5,11.5,16}
  	{\OpCodeGrpMEM,\OpCodeSUB, RegR, 1, RegB, dw, RegX, 0}

	\end{flushleft}  
}

\instructionentry{Description}{

The \texttt{SUB} instruction subtracts a signed 64-bit operand from a general register \texttt{RegB} and stores the result to the target register \texttt{RegR}. The instruction supports the immediate, the register and the two memory reference instruction formats. An arithmetic overflow results in a numeric overflow trap. The indexed formats may cause a memory reference associated trap.
}

\instructionentry{Operation}{

  	\texttt{RegR <- RegB - RegA} {(register format)}\\
  	\smallskip
  	\texttt{RegR <- RegB - immOperand( )} {(immediate format)}\\
   	\smallskip
  	\texttt{RegR <- RegR - memOperand( )} {(indexed formats)}\\
}

\instructionentry{Exceptions}{

 	Integer Overflow Trap\\
	Data Alignment Trap\\
	Memory Access Violation Trap\\
	Memory Protection Violation Trap\\
	Data TLB Miss Trap
}

\instructionentry{Notes}{

	None. 
}




